\chapter{DDL - Data Definition Language}
\section{Notazione}
\begin{itemize}
    \item \textbf{TERMINI DEL LINGUAGGIO} in maiuscolo
    \item \texttt{<x>} usato per isolare un termine x
    \item \texttt{[x]} indicano che il termine x è opzionale
    \item \texttt{\{x\}} indicano che il termine x può essere ripetuto da 0 a un numero arbitrario di volte
    \item \texttt{|} separa opzioni alternative
\end{itemize}

\section{Schemi}
\subsection{Cosa sono?}
Uno schema di base di dati è una collezione di oggetti: domini, tabelle, asserzioni, viste, privilegi, ...

\subsection{Sintassi di creazione}
\begin{lstlisting}[language=sql_generic]
CREATE SCHEMA
    [NomeSchema] -- Se omesso e' il nome dell'utente che ha lanciato il comando
    [[authorization] Autorizzazione] -- Termine del linguaggio
    {DefinizioneElementoSchema} -- Termine variabile
\end{lstlisting}

Dopo il comando \textbf{CREATE SCHEMA} compaiono le definizioni dei vari elementi.

Non è necessario che la definizione di tutti gli elementi avvenga contemporaneamente alla creazione dello schema, può avvenire in più fasi successive

\subsection{Esempio: Creazione di uno schema vuoto}
\begin{lstlisting}[language=sql_generic]
    CREATE SCHEMA GestioneAzienda;
\end{lstlisting}

\subsection{Esempio: Creazione di uno schema utilizzando tutte le opzioni disponibili}
\begin{lstlisting}[language=sql_generic]
CREATE SCHEMA GestioneAzienda -- Nome dello schema
AUTHORIZATION utente_admin -- Autorizzazione
CREATE TABLE Impiegato ( -- Termine variabile
    id_impiegato INT PRIMARY KEY,
    nome VARCHAR(50),
    cognome VARCHAR(50)
)
CREATE VIEW VistaImpiegati AS -- Altro termine variabile
SELECT nome, cognome FROM Impiegato;
\end{lstlisting}

\section{Tabelle}
\subsection{Cosa sono?}
Una tabella è costituita da una collezione ordinata di attributi e da un insieme (eventualmente vuoto) di vincoli
\subsection{Sintassi di creazione}
\begin{lstlisting}[language=sql_generic]
CREATE TABLE NomeTabella(
    NomeAttributo DOMINIO [ValoreDiDefault][Vincoli]
    {NomeAttributo DOMINIO [ValoreDiDefault][Vincoli]}
    [AltriVincoli]
)
\end{lstlisting}

\subsection{Esempio: Creazione di una tabella con elementi minimi}
\begin{lstlisting}[language=sql_generic]
CREATE TABLE Studente (
    matricola INT PRIMARY KEY,
    nome VARCHAR(30)
);
\end{lstlisting}
\subsection{Esempio: Creazione di una tabella full optional}
\begin{lstlisting}[language=sql_generic]
CREATE TABLE Esame (
    codice_esame INT,
    nome_corso VARCHAR(50) NOT NULL,
    data_esame DATE DEFAULT CURRENT_DATE,
    voto INT CHECK (voto BETWEEN 18 AND 30),
    matricola_studente INT,
    
    -- Definizione della chiave primaria composta
    CONSTRAINT PK_Esame PRIMARY KEY (codice_esame, matricola_studente),
    
    -- Definizione della chiave esterna con vincolo di integrita referenziale
    CONSTRAINT FK_Studente FOREIGN KEY (matricola_studente) 
        REFERENCES Studente(matricola)
        ON DELETE CASCADE
        ON UPDATE CASCADE
);
\end{lstlisting}

\section{Definizione dei Dati: I Domini}
\subsection{Cosa sono?}
I domini specificano i valori ammessi di ciascun attributo
\subsection{Domini elementari}
\begin{itemize}
    \item \textbf{Testuali}
    \begin{itemize}
        \item \textbf{varchar(n):} Stringa di lunghezza variabile tra 0 e n caratteri
        \item \textbf{char(n):} Stringa di lunghezza pari ad n
    \end{itemize}
    \item \textbf{Numerici}
    \begin{itemize}
        \item \textbf{Numerici Esatti}
        \begin{itemize}
            \item \textbf{integer / smallint :} rappresentano valori interi
            \item \textbf{numeric / decimal :} rappresentano i valori decimali
        \end{itemize}
        \item \textbf{Numerici approssimati}
        \begin{itemize}
            \item \textbf{real / double precision :} rappresentano valori a singola / doppia precisione in virgola mobile
            \item \textbf{float} permette di richiedere la precisione che si desidera
        \end{itemize}
    \end{itemize}
    \item \textbf{boolean:} un booleano
    \item \textbf{blob, clob:} oggetti di grandi dimensioni
    \begin{itemize}
        \item costruiti da binari (blob)
        \item costruiti da caratteri (clob)
    \end{itemize}
    vengono trattati e memorizzati diversamente dagli altri, vengono usati per informazioni non strutturate o multimediali
\end{itemize}

\section{Domini definiti dagli utenti}
\subsection{Cosa sono?}
Sono dei domini definiti dall'utente partendo dai domini elementari

\subsection{Sintassi di creazione}
\begin{lstlisting}[language=sql_generic]
CREATE DOMAIN NomeDominio AS DominioElementare
    [ValoreDefault]
    [Constraints]    
\end{lstlisting}

\subsection{Esempio: Creazione di dominio semplice}
\begin{lstlisting}[language=sql_generic]
CREATE DOMAIN Voto AS INT;
\end{lstlisting}
\subsection{Esempio: Creazione di un dominio full optional}
\begin{lstlisting}[language=sql_generic]
CREATE DOMAIN Voto AS SMALLINT
DEFAULT '0'
CHECK (value >= 18 AND value <= 30);
\end{lstlisting}

\section{Vincoli intrarelazionale}
Un vincolo intrarelazione è una regola che specifica delle condizioni all'interno di una relazione

\subsection{Tipidi di vincoli intrarelazionali}
Proprietà sempre valida all'interno di una relazione
\begin{itemize}
    \item \textbf{NOT NULL:} Il valore deve essere non nullo
    \item \textbf{UNIQUE:} I valori devono essere non ripetuti
    \item \textbf{PRIMARY KEY:} Chiave primaria (una sola per tabella)
    \item \textbf{CHECK:} Definisce condizioni complesse (sia intra che inter relazionali)
\end{itemize}

\subsection{Esempio}
\begin{lstlisting}[language=sql_generic]
CREATE TABLE Dipendenti (
    -- 1. VINCOLO DI CHIAVE PRIMARIA (e di dominio NOT NULL implicito)
    ID_Dipendente INT PRIMARY KEY,
    
    -- 2. VINCOLI DI DOMINIO (Tipo dato, NOT NULL e lunghezza)
    Nome VARCHAR(50) NOT NULL,
    Cognome VARCHAR(50) NOT NULL,
    
    -- 3. VINCOLO DI CHIAVE UNICA (Unique)
    -- L'email deve essere unica per ogni dipendente, ma potrebbe essere temporaneamente assente (NULL)
    Email VARCHAR(100) UNIQUE,
    
    -- 4. VINCOLO DI DOMINIO con controllo sul range (CHECK)
    -- Lo stipendio deve essere un numero positivo
    Stipendio_Mensile DECIMAL(10, 2) CHECK (Stipendio_Mensile > 0),
    
    Data_Assunzione DATE NOT NULL,
    Data_Fine_Contratto DATE,
    
    -- 5. VINCOLO DI TUPLA (CHECK che coinvolge piu colonne)
    -- La data di fine contratto deve essere successiva alla data di assunzione
    CONSTRAINT CHK_Date_Contratto CHECK (Data_Fine_Contratto > Data_Assunzione)
);
\end{lstlisting}

\newpage
\section{Vincoli interrelazionali: Integrità referenziale}
\subsection{Cosa sono?}
Esprimono un legame gerarchico (padre / figlio) fra tabelle.

Alcuni attributi della tabella figlio sono definti \textbf{FOREIGN KEY} si devono riferire (\textbf{REFERENCES}) ad alcuni attributi della tabella padre che costituiscono una chiave.

I valori contenuti nella \textbf{FOREIGN KEY} devono essere sempre presenti nella tabella padre.

\subsection{Sintassi}
\subsubsection{REFERENCES}
Nella parte di definizione degli attributi con il costrutto sintattico \textbf{REFERENCES}
\begin{lstlisting}[language=sql_generic]
AttrFiglio CHAR(3) REFERENCES TabellaPadre(AttrPadre)
\end{lstlisting}

\subsubsection{FOREIGN KEY}
Oppure dopo la definizione degli attributi con i costruttori \textbf{FOREGIN KEY} e \textbf{REFERENCES}
\begin{lstlisting}[language=sql_generic]
FOREIGN KEY(AttrFiglio)
    REFERENCES TabellaPadre(AttrPadre)
\end{lstlisting}

Quando si hanno più attributi da riferire, si utilizza sempre \textbf{FOREIGN KEY}
\begin{lstlisting}[language=sql_generic]
FOREIGN KEY(AttrFiglio1{, AttrFiglio2})
    REFERENCES TabellaPadre(AttrPadre1{, AttrPadre2})
\end{lstlisting}

Se si omettono gli attributi destinazione, vengono assunti quelli della chiave primaria
\begin{lstlisting}[language=sql_generic]
AttrFiglio CHAR(3) REFERENCES TabellaPadre
\end{lstlisting}

\subsection{Esempio}
\begin{lstlisting}[language=sql_generic]
CREATE TABLE Infrazioni (
    Codice INTEGER PRIMARY KEY, -- Identificativo dell'infrazione (es. 34321)
    Data DATE NOT NULL,
    Vigile INTEGER NOT NULL,
    Prov CHAR(2),
    Numero CHAR(6),
    
    -- Vincolo interrelazionale singolo (in linea o a livello tabella)
    FOREIGN KEY (Vigile) REFERENCES Vigili(Matricola),
    
    -- Vincolo interrelazionale composto (obbligatoriamente a fine tabella)
    FOREIGN KEY (Prov, Numero) REFERENCES Auto(Prov, Numero)
);
\end{lstlisting}


\section{Aggiornamento}
\subsection{Il problema dell'aggiornamento}
\begin{figure}[htbp]
    \centering
    \begin{tikzpicture}[remember picture, >=Stealth, thick]

        % Tabella: Studente
        \node (studente) [label={below:{\Large\sffamily Studente}}] {
            \sffamily\color{textdarkblue}
            \renewcommand{\arraystretch}{1.3}
            \begin{tabular}{|C{1.2cm}|C{1.6cm}|C{1.2cm}|C{1.0cm}|}
                \hline
                \rowcolor{headerblue} \underline{\textbf{Matr}} & \textbf{Nome} & \textbf{Città} & \textbf{CDip} \\ \hline
                \rowcolor{rowblue} 34321 & Luca & Mi & Inf \\ \hline
                \rowcolor{rowblue} 53524 & Giovanni & To & Mat \\ \hline
                \rowcolor{rowblue} 64521 & Emilio & Ge & Ing \\ \hline
                \rowcolor{rowblue} 73321 & Francesca & Vr & Mat \\ \hline
            \end{tabular}
        };

        % Tabella: Esame
        \node (esame) [right=2cm of studente, label={above:{\Large\sffamily Esame}}] {
            \sffamily\color{textdarkblue}
            \renewcommand{\arraystretch}{1.3}
            \begin{tabular}{|C{1.2cm}|C{1.8cm}|C{1.6cm}|C{1.0cm}|}
                \hline
                \rowcolor{headerblue} \underline{\textbf{Matr}} & \underline{\textbf{CodCorso}} & \textbf{Data} & \textbf{Voto} \\ \hline
                \rowcolor{rowblue} 34321 & 1 & 25/02/01 & 28 \\ \hline
                \rowcolor{rowblue} 34321 & 2 & 14/12/01 & 30 \\ \hline
                \rowcolor{rowblue} 64521 & 1 & 12/03/02 & 24 \\ \hline
                \rowcolor{rowblue} 73321 & 4 & 18/01/02 & 18 \\ \hline
            \end{tabular}
        };

        % Freccia rossa per la Foreign Key
        \draw[red, ->, line width=1.5pt] ([xshift=0.6cm, yshift=1.1cm]esame.west) -- ([xshift=-1.3cm, yshift=1.5cm]esame.west);

    \end{tikzpicture}
\end{figure}

Se elimino da Studente la tupla con matricola 34321, cosa succede alla tabella esame?

\subsection{Il problema delle violazioni}
Se viene eseguita una operazione di aggiornamento che porta alla violazione di un vincolo di integrità referenziale la base di dati diventa non valida.

Per evitare ciò, vengono definite un insieme di politiche per evitare che questo accada, permettendo di scegliere delle reazioni da adottare in caso di violazione.

Queste politiche vanno dichiarate nella dichiarazione DDL dello schema.

Per tutti gli altri tipoli di vicnoli invece a seguito di una violazione il comando di aggiornamento viene rifiutato segnalando l'errore.

\vspace{\baselineskip}

\begin{figure}[htbp]
    \centering
    \sffamily
    
    % Colonna di sinistra (43% della larghezza)
    \begin{minipage}[c]{0.43\textwidth}
        \textbf{Tipo di violazione}\\
        Inserimento o modifica di tupla che viola il vincolo sulla \textcolor{red!80!black}{\textbf{tabella figlio}} (interna)
    \end{minipage}%
    %
    % Spazio centrale con la freccia (10% della larghezza)
    \begin{minipage}[c]{0.10\textwidth}
        \centering
        \begin{tikzpicture}
            \draw[-{Stealth}, red!80!black, line width=1.2pt] (0,0) -- (0.8,0);
        \end{tikzpicture}
    \end{minipage}%
    %
    % Colonna di destra (43% della larghezza)
    \begin{minipage}[c]{0.43\textwidth}
        \textbf{Politica di reazione}\\
        In entrambi i casi il comando non viene eseguito
    \end{minipage}
\end{figure}
\begin{figure}[htbp]
    \centering
    \sffamily
    
    % Colonna di sinistra (43% della larghezza)
    \begin{minipage}[c]{0.43\textwidth}
        Inserimento o modifica di tupla che viola il vincolo sulla tabella padre (esterna)
    \end{minipage}%
    %
    % Spazio centrale con la freccia (10% della larghezza)
    \begin{minipage}[c]{0.10\textwidth}
        \centering
        \begin{tikzpicture}
            \draw[-{Stealth}, red!80!black, line width=1.2pt] (0,0) -- (0.8,0);
        \end{tikzpicture}
    \end{minipage}%
    %
    % Colonna di destra (43% della larghezza)
    \begin{minipage}[c]{0.43\textwidth}
        Quattro politiche possibili
    \end{minipage}
\end{figure}

\subsection{Reazioni a violazioni sulla tabelal padre}
\begin{itemize}
    \item \textbf{CASCADE:} L'operazione di modifica / cancellazione viene propagata alla tabella intera
    \item \textbf{SET NULL:} NULL nella tabella figlio in entrambi i casi
    \item \textbf{SET DEFAULT:} Default nella tabella figlio in entrambi i casi
    \item \textbf{NO ACTION:} Rifiutata in entrambi i casi
\end{itemize}

\subsection{Vincoli con nomi}
A fini diagnostici è spesso utilie sapere quale vincolo è stato violato a seguito di un'azione sul DB.

A tale scopo è possibile associare dei nomi ai vincoli, ad esempio:
\begin{lstlisting}[language=sql_generic]
Stipendio INTEGER CONSTRAINT StipendioPositivo
CHECK (Stipendio > 0)
\end{lstlisting}
\begin{lstlisting}[language=sql_generic]
CONSTRAINT ForeignKeySedi
FOREIGN KEY (Sede) REFERENCES Sedi
\end{lstlisting}

\section{Modifiche degli schemi: ALTER}
\subsection{A cosa serve?}
Modifica oggetti

\subsection{Tipi di ALTER}
\begin{itemize}
    \item \textbf{ALTER DOMAIN:} Modifica domini
    \item \textbf{ALTER TABLE:} Modifica tabelle
\end{itemize}

Quando si definisce un nuovo vincolo deve essere soddisfatto dalla istanza presente, altrimenti viene rifiutato

\subsection{Sintassi ALTER DOMAIN}
\begin{lstlisting}[language=sql_generic]
ALTER DOMAIN NomeDominio <
    SET      DEFAULT ValoreDefault |
    DROP     DEFAULT |
    ADD      CONSTRAINT DefVincolo |
    DROP     CONSTRAINT NomeVincolo >
\end{lstlisting}

\subsection{Esempio di ALTER DOMAIN}
\begin{lstlisting}[language=sql_generic]
-- 1. Definiamo un valore di default standard per tutti i nuovi inserimenti
ALTER DOMAIN VotoEsame 
    SET DEFAULT 18;
-- 2. Aggiungiamo un vincolo per gestire anche l'estensione del voto con la "Lode" (es. inserendo 31)
ALTER DOMAIN VotoEsame 
    ADD CONSTRAINT ControlloLode CHECK (VALUE >= 18 AND VALUE <= 31);
-- 3. Rimuoviamo il valore di default inserito in precedenza se non e piu necessario
ALTER DOMAIN VotoEsame 
    DROP DEFAULT;
-- 4. Rimuoviamo il vincolo sulla lode per tornare alle impostazioni base
ALTER DOMAIN VotoEsame 
    DROP CONSTRAINT ControlloLode;
\end{lstlisting}

\subsection{Sintassi ALTER TABLE}
\begin{lstlisting}[language=sql_generic]
ALTER TABLE NomeTabella <
    ALTER COLUMN NomeAttributo <
        SET DEFAULT NuovoDefault |
        DROP DEFAULT > |
    DROP COLUMN NomeAttributo |
    ADD COLUMN DefAttributo |
    DROP CONSTRAINT NomeVincolo |
    ADD CONSTRAINT DefVincolo >
\end{lstlisting}

\subsection{Esempio di ALTER TABLE}
\begin{lstlisting}[language=sql_generic]
-- 1. Aggiungiamo una nuova colonna per memorizzare l'Email dello studente
ALTER TABLE Studente 
    ADD COLUMN Email VARCHAR(100);
-- 2. Modifichiamo la colonna "Citta" esistente per impostare un valore di default predefinito
ALTER TABLE Studente 
    ALTER COLUMN Citta SET DEFAULT 'Mi';
-- 3. Rimuoviamo la colonna "CDip" (Codice Dipartimento) se la tabella viene ristrutturata
ALTER TABLE Studente 
    DROP COLUMN CDip;
-- 4. Aggiungiamo un vincolo di unicita (Unique) sulla colonna Email appena creata
ALTER TABLE Studente 
    ADD CONSTRAINT EmailUnica UNIQUE (Email);
-- 5. Rimuoviamo un vecchio vincolo di controllo (se esistente) identificato dal suo nome
ALTER TABLE Studente 
    DROP CONSTRAINT VerificaMatricola;
\end{lstlisting}

\section{Vincolo interrelazionale: UPDATE}

Quando si effettua una modifica (\texttt{UPDATE}) sui dati all'interno di una base di dati relazionale, è fondamentale garantire che i legami tra le tabelle rimangano consistenti. Il vincolo di integrità referenziale (Foreign Key) interviene in modo specifico per regolamentare cosa succede quando viene modificata una chiave.

\subsection{A cosa serve}
Il vincolo interrelazionale durante un'operazione di modifica ha lo scopo di impedire la creazione di ``tuple orfane''. Si considerano due scenari distinti:
\begin{itemize}
    \item \textbf{Modifica sulla tabella figlio (interna):} Se si modifica il valore della chiave esterna nella tabella figlio, il database controlla semplicemente che il nuovo valore esista nella chiave primaria della tabella padre. Se non esiste, l'operazione viene bloccata.
    \item \textbf{Modifica sulla tabella padre (esterna):} Se si modifica il valore della chiave primaria nella tabella padre, tutti i record della tabella figlio che puntavano a quel vecchio valore rimarrebbero associati a una chiave non più esistente. Per evitare questa violazione, il DBMS applica delle specifiche \textit{politiche di reazione}.
\end{itemize}

\subsubsection{Sintassi dell'operazione}
La definizione della politica di reazione per la modifica viene dichiarata in fase di creazione o alterazione della tabella figlio, subito dopo la specifica della chiave esterna (\texttt{FOREIGN KEY}).

La sintassi generica all'interno del linguaggio SQL DDL è la seguente:
\begin{lstlisting}[language=sql_generic]
FOREIGN KEY (AttributoTabellaFiglio) REFERENCES TabellaPadre(AttributoTabellaPadre) ON UPDATE [ RESTRICT | CASCADE | SET NULL | SET DEFAULT ]
\end{lstlisting}

Il significato delle quattro politiche possibili è il seguente:
\begin{enumerate}
    \item \textbf{RESTRICT (o NO ACTION):} È il comportamento di default. La modifica sulla tabella padre viene rifiutata e solleva un errore se esistono tuple collegate nella tabella figlio.
    \item \textbf{CASCADE:} La modifica viene propagata a cascata. Modificando il valore nella tabella padre, tutti i record corrispondenti nella tabella figlio verranno aggiornati automaticamente con il nuovo valore.
    \item \textbf{SET NULL:} Il valore della chiave esterna nei record correlati della tabella figlio viene impostato a \texttt{NULL} (richiede che la colonna non abbia il vincolo \texttt{NOT NULL}).
    \item \textbf{SET DEFAULT:} Il valore della chiave esterna nella tabella figlio viene reimpostato al valore di default precedentemente configurato per quella colonna.
\end{enumerate}

\subsection{Esempi pratici}

Riprendiamo come riferimento le tabelle \textbf{Studente} (tabella padre, con chiave primaria \texttt{Matr}) ed \textbf{Esame} (tabella figlio, con chiave esterna \texttt{Matr} che fa riferimento a \texttt{Studente}).

\subsubsection{Esempio 1: Politica CASCADE (Propagazione automatica)}
Si supponga di voler modificare la matricola dello studente `Luca' da \texttt{34321} a \texttt{99999}. Configurando la tabella con la politica \texttt{CASCADE}:
\begin{lstlisting}[language=sql_generic]
ALTER TABLE Esame 
    ADD CONSTRAINT fk_studente 
    FOREIGN KEY (Matr) REFERENCES Studente(Matr)
    ON UPDATE CASCADE;
\end{lstlisting}
Eseguendo l'istruzione \texttt{UPDATE Studente SET Matr = 99999 WHERE Matr = 34321;}, il database aggiornerà automaticamente anche tutte le righe all'interno della tabella \texttt{Esame} in cui la matricola era \texttt{34321}, sostituendola con \texttt{99999}.

\medskip
\subsubsection{Esempio 2: Politica RESTRICT (Blocco cautelativo)}
Nel caso in cui le matricole debbano rimanere tassativamente immutabili per ragioni storiche e di sicurezza, si opta per \texttt{RESTRICT}:
\begin{lstlisting}[language=sql_generic]
ALTER TABLE Esame 
    ADD CONSTRAINT fk_studente 
    FOREIGN KEY (Matr) REFERENCES Studente(Matr)
    ON UPDATE RESTRICT;
\end{lstlisting}
Se si provasse ad eseguire la modifica della matricola \texttt{34321} nella tabella \texttt{Studente}, il DBMS bloccherebbe l'operazione restituendo un errore del tipo: \textit{``Violazione del vincolo di chiave esterna: record correlati presenti nella tabella Esame''}.

\subsubsection{Esempio 3: Politica SET NULL (Annullamento del legame)}
Se l'esame sostenuto deve rimanere registrato nel sistema ma, in seguito alla modifica o riorganizzazione di una targa/codice identificativo padre, si decidesse temporaneamente di scollegarlo dallo studente:
\begin{lstlisting}[language=sql_generic]
ALTER TABLE Esame 
    ADD CONSTRAINT fk_studente 
    FOREIGN KEY (Matr) REFERENCES Studente(Matr)
    ON UPDATE SET NULL;
\end{lstlisting}
Cambiando la matricola nella tabella padre, le righe corrispondenti nella tabella \texttt{Esame} vedranno il proprio campo \texttt{Matr} trasformarsi in \texttt{NULL}, indicando che l'esame è al momento ``orfano'' di uno studente identificato.

\section{Vincolo interrelazionale: DROP}
\subsection{Cosa fa?}
Cancella oggetti dallo schema

Ha 2 opzioni:
\begin{itemize}
    \item \textbf{RESTRICT:} Un oggetto non è rimosso se non è vuoto
    \item \textbf{CASCADE:} Viene rimosso l'oggetto e tutto ciò che è coinvolto nella definizione dell'oggetto
\end{itemize}

\subsection{DROP DOMAIN}
Cancellazione del dominio

\subsubsection{Sintassi}
\begin{lstlisting}[language=sql_generic]
DROP DOMAIN NomeDominio [ RESTRICT | CASCADE ];
\end{lstlisting}

\subsubsection{Esempio di utilizzo}
\begin{lstlisting}[language=sql_generic]
-- 1. Creazione del dominio
CREATE DOMAIN CodiceDipartimento AS CHAR(3)
    CHECK (VALUE IS NOT NULL);

-- 2. Utilizzo del dominio nella tabella "Studente"
CREATE TABLE Studente (
    Matr INT PRIMARY KEY,
    Nome VARCHAR(50),
    CDip CodiceDipartimento -- La colonna CDip usa il dominio appena creato
);

DROP DOMAIN CodiceDipartimento RESTRICT;
\end{lstlisting}

\subsubsection{DROP TABLE}
Cancellazione della tabella

\subsubsection{Sintassi}
\begin{lstlisting}
DROP TABLE NomeTabella [ RESTRICT | CASCADE ];
\end{lstlisting}

\subsubsection{Esempi}
\begin{lstlisting}[language=sql_generic]
-- Creazione della tabella Padre
CREATE TABLE Studente (
    Matr INT PRIMARY KEY,
    Nome VARCHAR(50)
);

-- Creazione della tabella Figlio con il vincolo di Foreign Key
CREATE TABLE Esame (
    CodCorso INT,
    Matr INT,
    Voto INT,
    PRIMARY KEY (CodCorso, Matr),
    CONSTRAINT fk_studente FOREIGN KEY (Matr) REFERENCES Studente(Matr)
);

DROP TABLE Studente RESTRICT;
-- Nota: Scrivere "DROP TABLE Studente;" produce lo stesso effetto
\end{lstlisting}